\documentclass{article}
\usepackage[margin=25mm]{geometry}
\usepackage{booktabs}
\usepackage{amsmath}
\title{Three-cloud propagation prototype: concurrency benchmark}
\author{NAUSICAA analysis}
\date{29 October 2026}
\begin{document}
\maketitle
\section*{Method}
For each of 100 events from each 1~mm source sample (double beta and single
electrons), the Julia prototype produced an ion cloud and electron clouds using
the pure-Xe and Xe--15\%~He Petit diffusion coefficients. The event location was
drawn once per source event; ion and electron distances sum to 150~cm. The
benchmark includes Julia process startup, HDF5 source reads, diffusion and
compressed HDF5 output. Each stream writes a separate file.

\section*{Results}
\begin{center}
\begin{tabular}{rrrrr}
\toprule
Workers & Events/stream & Wall time (s) & Speedup & Output (MB) \\
\midrule
1 & 100 & 39.03 & 1.00 & 64.2 \\
2 & 100 & 19.98 & 1.95 & 64.2 \\
4 & 100 & 13.59 & 2.87 & 64.2 \\
6 & 100 & 7.63 & 5.12 & 64.2 \\
\bottomrule
\end{tabular}
\end{center}
The sum of individual peak resident-memory measurements for six processes is 4.90 GB. This sum is a conservative ceiling for simultaneous resident memory, not a measured concurrent peak.

\section*{Validation and limits}
All 100 events in each source sample have the same sampled ion drift distance
across their three clouds. Ion and electron drift distances sum to 150~cm.
Every cloud conserves its sampled carriers across retained and outside-bin
counts. Per-event arrays and counts match exactly for one, two, four and six
workers. These measurements cover 100 events per stream and include output
to the local filesystem; production-scale throughput and concurrent peak
resident memory require a separate full-run measurement.
\end{document}
